\textbf{Keywords:} IP, xschem, ngspice, Corners, Layout, DRC, LVS, LPE,
Documentation

Before you start the tutorial you need to have the tools installed.

If you're a student of mine, then you must install the tools locally on
your PC. Read check \href{https://analogicus.com/aic2026/the_tools}{The
Tools} chapter first.

If you're the impatient kind, then check
\url{https://analogicus.com/aicex/} which shows how to start the docker
image.

\section{Create the IP}\label{create-the-ip}\index{create the ip@Create the IP}

I've made some scripts to automatically generate the IP.

To see what files are generated, see
\texttt{tech\_sky130A/cicconf/lelo.yaml}

\begin{Shaded}
\begin{Highlighting}[]
\BuiltInTok{cd}\NormalTok{ aicex/ip}
\ExtensionTok{cicconf}\NormalTok{ newip ex }\AttributeTok{{-}{-}project}\NormalTok{ lelo }\AttributeTok{{-}{-}technology}\NormalTok{ sky130A }\AttributeTok{{-}{-}ip}\NormalTok{  tech\_sky130A/cicconf/lelo.yaml}
\end{Highlighting}
\end{Shaded}

\section{The file structure}\label{the-file-structure}\index{the file structure@The file structure}

It matters how you name files, and store files. I would be surprised if
you had a good method already, as such, I won't allow you to make your
own folder structure and names for things. I also control the filenames
and folder structure because there are many scripts to make your life
easier (yes, really) that rely on an exact structure. Don't mess with
it.

\subsection{Github workflows}\label{github-workflows}\index{github workflows@Github workflows}

On github it's possible use something called workflows to run things
every time you push a new version. It's really nice, since it can then
check that your design is valid.

The workflows are defined below.

\begin{Shaded}
\begin{Highlighting}[]
\ExtensionTok{.github}
   \ExtensionTok{workflows}  
   \ExtensionTok{docs.yaml} \CommentTok{\# Generate a github page }
   \ExtensionTok{drc.yaml}  \CommentTok{\# Run Design Rule Checks }
   \ExtensionTok{gds.yaml}  \CommentTok{\# Generate a GDS file from layout }
   \ExtensionTok{lvs.yaml}  \CommentTok{\# Run Layout Versus Schematic }
             \CommentTok{\# and Layout Parasitic Extraction}
\end{Highlighting}
\end{Shaded}

\subsection{Configuration files}\label{configuration-files}\index{configuration files@Configuration files}

Each IP has a few files that define the setup, you'll need to modify at
least the \texttt{README.md} and the \texttt{info.yaml}.

\begin{Shaded}
\begin{Highlighting}[]
 \ExtensionTok{.gitignore}  \CommentTok{\# files that are ignored by git}
 \ExtensionTok{README.md}   \CommentTok{\# Frontpage documentation }
 \ExtensionTok{config.yaml} \CommentTok{\# What libraries are used. Used by  cicconf}
 \ExtensionTok{info.yaml}   \CommentTok{\# Setup names, authors etc }
 \ExtensionTok{media}       \CommentTok{\# Where you should store images for documentation}
 \ExtensionTok{tech} \AttributeTok{{-}}\OperatorTok{\textgreater{}}\NormalTok{ ../tech\_sky130A  }\CommentTok{\# The technology library}
\end{Highlighting}
\end{Shaded}

\subsection{Design files}\label{design-files}\index{design files@Design files}

A ``cell'' in the open source EDA world should consist of the following
files

\begin{itemize}
\tightlist
\item
  Schematic (.sch)
\item
  Layout (.mag)
\item
  Documentation (.md)
\end{itemize}

The files must have the same name, and must be stored in
\texttt{design/\textless{}LIB\textgreater{}/} as shown below.

Note there are also two symbolic links to other libraries. These two
libraries contain standard cells and standard analog transistors (ATR)
that you should be using.

\begin{Shaded}
\begin{Highlighting}[]
\ExtensionTok{design}
  \ExtensionTok{LELO\_EX\_SKY130A}
  \ExtensionTok{LELO\_EX.sch}
  \ExtensionTok{JNW\_ATR\_SKY130A} \AttributeTok{{-}}\OperatorTok{\textgreater{}}\NormalTok{ ../../jnw\_atr\_sky130a/design/JNW\_ATR\_SKY130A}
  \ExtensionTok{JNW\_TR\_SKY130A} \AttributeTok{{-}}\OperatorTok{\textgreater{}}\NormalTok{ ../../jnw\_tr\_sky130a/design/JNW\_TR\_SKY130A}
\end{Highlighting}
\end{Shaded}

For example, if the cell name was \texttt{LELO\_EX}, then you would have

\begin{itemize}
\tightlist
\item
  \texttt{design/LELO\_EX\_SKY130A/LELO\_EX.sch}: Schematic (xschem)
\item
  \texttt{design/LELO\_EX\_SKY130A/LELO\_EX.sym}: Symbol (xschem)
\item
  \texttt{design/LELO\_EX\_SKY130A/LELO\_EX.mag}: Layout (Magic)
\item
  \texttt{design/LELO\_EX\_SKY130A/LELO\_EX.md} : Markdown documentation
  (any text editor)
\end{itemize}

All these files are text files, so you can edit them in a text editor,
but mostly you shouldn't (except for the Markdown)

\subsection{Simulations}\label{simulations}\index{simulations@Simulations}

All simulations shall be stored in \texttt{sim}. Once you have a
Schematic ready for simulation, then

\begin{Shaded}
\begin{Highlighting}[]
\BuiltInTok{cd}\NormalTok{ sim }
\FunctionTok{make}\NormalTok{ cell CELL=LELO\_EX}
\end{Highlighting}
\end{Shaded}

This will make a simulation folder for you. Repeat for all your cells.

\begin{Shaded}
\begin{Highlighting}[]
\ExtensionTok{sim}
  \ExtensionTok{Makefile}
  \ExtensionTok{cicsim.yaml} \AttributeTok{{-}}\OperatorTok{\textgreater{}}\NormalTok{ ../tech/cicsim/cicsim.yaml}
\end{Highlighting}
\end{Shaded}

\subsection{The work}\label{the-work}\index{the work@The work}

All commands (except for simulation), shall be run in the \texttt{work}
folder.

In the \texttt{work/} folder there are startup files for Xschem
(xschemrc) and Magic (.magicrc). They tell the tools where to find the
process design kit, symbols, etc. At some point you probably need to
learn those also, but I'd wait until you feel a bit more comfortable.

\begin{Shaded}
\begin{Highlighting}[]
\ExtensionTok{work}
 \ExtensionTok{.magicrc}
 \ExtensionTok{Makefile}
 \ExtensionTok{mos.24bit.dstyle} \AttributeTok{{-}}\OperatorTok{\textgreater{}}\NormalTok{ ../tech/magic/mos.24bit.dstyle}
 \ExtensionTok{mos.24bit.std.cmap} \AttributeTok{{-}}\OperatorTok{\textgreater{}}\NormalTok{ ../tech/magic/mos.24bit.std.cmap}
 \ExtensionTok{xschemrc}
\end{Highlighting}
\end{Shaded}

\section{Github setup}\label{github-setup}\index{github setup@Github setup}

Create a repository on \href{https://github.com}{github}. The name of
the repository that you make on GitHub has to be the same as what is
written after \texttt{\textless{}your\ username\textgreater{}} in the
last command below. In this example, that is \texttt{lelo\_ex\_sky130a}.

\begin{Shaded}
\begin{Highlighting}[]
\BuiltInTok{cd}\NormalTok{ lelo\_ex\_sky130a}
\FunctionTok{git}\NormalTok{ remote add origin }\DataTypeTok{\textbackslash{}}
\NormalTok{ git@github.com:}\OperatorTok{\textless{}}\NormalTok{your username}\OperatorTok{\textgreater{}}\NormalTok{/lelo\_ex\_sky130a.git}
\end{Highlighting}
\end{Shaded}

\section{Start working}\label{start-working}\index{start working@Start working}

\subsection{Edit README.md}\label{edit-readme.md}\index{edit readme.md@Edit readme.md}

Open README.md in your favorite text editor and make necessary changes.

\subsection{Familiarize yourself with the Makefile and
make}\label{familiarize-yourself-with-the-makefile-and-make}

I write all commands I do into a Makefile. There is nothing special with
a Makefile, it's just what I choose to use 20 years ago. I'm not sure
I'd choose something different now.

\begin{Shaded}
\begin{Highlighting}[]
\BuiltInTok{cd}\NormalTok{ work}
\FunctionTok{make}
\end{Highlighting}
\end{Shaded}

Take a look inside the file called Makefile.

\section{\texorpdfstring{Draw Schematic
}{Draw Schematic }}\label{draw-schematic}

The block we'll make is a current mirror with a 1 to 4 scaling.

A schematic is how we describe the connectivity, and the types of
devices in an analog circuit. The open source schematic editor we will
use is XSchem.

Open the schematic:

\begin{Shaded}
\begin{Highlighting}[]
\ExtensionTok{xschem} \AttributeTok{{-}b}\NormalTok{ ../design/LELO\_EX\_SKY130A/LELO\_EX.sch }\KeywordTok{\&}
\end{Highlighting}
\end{Shaded}

\subsection{Add Ports}\label{add-ports}\index{add ports@Add ports}

Add IBPS\_5U and IBNS\_20U ports, the P and N in the name signifies what
transistor the current comes from. So IBPS must go into a diode
connected NMOS, and N will be our output, and go into a diode connected
PMOS somewhere else.

\subsection{Add transistors}\label{add-transistors}\index{add transistors@Add transistors}

Use `I' or `Shift+i' (note the letter case) to open the library manager.
Click the \texttt{lelo\_ex\_sky130A/design} path, then
\texttt{JNW\_ATR\_SKY130A} and select \texttt{JNWATR\_NCH\_4C5F0.sym}

The naming convention for these transistors is
\texttt{\textless{}number\ of\ contacts\ on\ drain/source\textgreater{}C\textless{}times\ minimum\ gate\ length\textgreater{}F},
so the number before the C is the width, and the number before/after the
F is the length. The absolute size does not matter for now. Just think
``4C5F0 is a 4 contact wide long transistor'', while a ``4C1F2 is a 4
contact wide, short transistor''.

Select the transistor and press `c' to copy it, while dragging, press
`shift-f' to flip the transistor so our current mirror looks nice.
`shift-r' rotates the transistor, but we don't want that now.

Place two transistors for the output transistor, as shown in the figure
below.

Press ESC to deselect everything

Select the input transistor, and change the name to `xo1'

Select the first output transistor, and change the name to
`xo0{[}1:0{]}'. Using bus notation on the name will create 2
transistors.

Select the second output transistor and give it the name `xo2{[}1:0{]}'.

Select ports, and use `m' to move the ports close to the transistors.

Press `w' to route wires.

Use `shift-z' and z, to zoom in and out

Use `f' to zoom full screen

Remember to save the schematic


{
\centering
\aicfig{media/LELO_EX.png}
}

\emph{Figure 1: Finished LELO\_EX schematic in Xschem - three NMOS
current mirror devices between the IBPS\_5U, IBNS\_20U and VSS ports}

\subsection{Netlist schematic}\label{netlist-schematic}\index{netlist schematic@Netlist schematic}

Check that the netlist looks OK

In work/

\begin{Shaded}
\begin{Highlighting}[]
\FunctionTok{make}\NormalTok{ xsch CELL=LELO\_EX}
\FunctionTok{cat}\NormalTok{ xsch/LELO\_EX.spice}
\end{Highlighting}
\end{Shaded}

\section{\texorpdfstring{Typical corner SPICE simulation
}{Typical corner SPICE simulation }}\label{typical-corner-spice-simulation}

I've made \href{https://github.com/wulffern/cicsim}{cicsim} that I use
to run simulations (ngspice) and extract results

\subsection{Setup simulation
environment}\label{setup-simulation-environment}

Navigate to the \texttt{lelo\_ex\_sky130a/sim/} directory.

Make a new simulation folder

\begin{Shaded}
\begin{Highlighting}[]
\ExtensionTok{cicsim}\NormalTok{ simcell  LELO\_EX\_SKY130A LELO\_EX }\DataTypeTok{\textbackslash{}}
\NormalTok{    ../tech/cicsim/cell\_spice/template.yaml}
\end{Highlighting}
\end{Shaded}

I would recommend you have a look at the template.yaml file to
understand what happens.

\subsection{Familiarize yourself with the simulation
folder}\label{familiarize-yourself-with-the-simulation-folder}

I've added quite a few options to cicsim, and it might be confusing. For
reference, these are what the files are used for

\begin{longtable}[]{@{}
  >{\raggedright\arraybackslash}p{(\linewidth - 2\tabcolsep) * \real{0.2154}}
  >{\raggedright\arraybackslash}p{(\linewidth - 2\tabcolsep) * \real{0.7846}}@{}}
\toprule\noalign{}
\begin{minipage}[b]{\linewidth}\raggedright
File
\end{minipage} & \begin{minipage}[b]{\linewidth}\raggedright
Description
\end{minipage} \\
\midrule\noalign{}
\endhead
\bottomrule\noalign{}
\endlastfoot
Makefile & Simulation commands \\
cicsim.yaml & Setup for cicsim \\
summary.yaml & Generate a README with simulation results \\
tran.meas & Measurement to be done after simulation \\
tran.py & Optional python script to run for each simulation \\
tran.spi & Transient testbench \\
tran.yaml & What measurements to summarize \\
\end{longtable}

The default setup should run, so

\begin{Shaded}
\begin{Highlighting}[]
\BuiltInTok{cd}\NormalTok{ LELO\_EX}
\FunctionTok{make}\NormalTok{ typical}
\end{Highlighting}
\end{Shaded}

\subsection{Modify default testbench
(tran.spi)}\label{modify-default-testbench-tran.spi}

Delete the VDD source

Add a current source of 5uA, and a voltage source of 1V to IBNS\_20U

\begin{Shaded}
\begin{Highlighting}[]
\NormalTok{IBP 0 IBPS\_5U dc 5u}
\NormalTok{V0  IBNS\_20U 0 dc 1}
\end{Highlighting}
\end{Shaded}

Save the current in V0 by adding i(V0) to the save statement in the
testbench

Save the voltage by adding v(IBPS\_5U) to the save statement

\begin{Shaded}
\begin{Highlighting}[]
\NormalTok{.save i(V0) v(IBPS\_5U)}
\end{Highlighting}
\end{Shaded}

\subsection{Modify measurements
(tran.meas)}\label{modify-measurements-tran.meas}

Add measurement of the current and VGS. It must be added between the
``MEAS\_START'' and ``MEAS\_END'' lines.

\begin{Shaded}
\begin{Highlighting}[]
\NormalTok{let ibn = {-}i(v0)}
\NormalTok{meas tran ibns\_20u find ibn at=5n}
\NormalTok{meas tran vgs\_m1 find v(ibps\_5u) at=5n}
\end{Highlighting}
\end{Shaded}

Run simulation

\begin{Shaded}
\begin{Highlighting}[]
\FunctionTok{make}\NormalTok{ typical}
\end{Highlighting}
\end{Shaded}

and check that the output looks okish.

Try to run the simulation again

\begin{Shaded}
\begin{Highlighting}[]
\FunctionTok{make}\NormalTok{ typical}
\end{Highlighting}
\end{Shaded}

If everything works, then the simulation now should \textbf{not} be run.
Every time cicsim runs (provided the \texttt{sha:\ True} option is set
in \texttt{cicsim.yaml}) cicsim will compute a SHA hash of all files
(stored in output\_tran/\emph{.sha}) that is referenced in the
\texttt{tran.spi}. Next time cicsim is run, it checks the hash's and
does not re-run if there is no need (no files changed).

Sometimes you want to force running, and you can do that by

\begin{Shaded}
\begin{Highlighting}[]
\FunctionTok{make}\NormalTok{ typical OPT=}\StringTok{"{-}{-}no{-}sha"}
\end{Highlighting}
\end{Shaded}

Often, it's the measurement that I get wrong, so instead of rerunning
simulation every time I've added a ``--no-run'' option to cicsim. For
example

\begin{Shaded}
\begin{Highlighting}[]
\FunctionTok{make}\NormalTok{ typical OPT=}\StringTok{"{-}{-}no{-}run"}
\end{Highlighting}
\end{Shaded}

will skip the simulation, and rerun only the measurement. This is why
you should split the testbench and the measurement. Simulations can run
for days, but measurement takes seconds.

\subsection{Modify result specification
(tran.yaml)}\label{modify-result-specification-tran.yaml}

Add the result specifications, for example

\begin{Shaded}
\begin{Highlighting}[]
\FunctionTok{ibn}\KeywordTok{:}
\AttributeTok{  }\FunctionTok{src}\KeywordTok{:}
\AttributeTok{    }\KeywordTok{{-}}\AttributeTok{ ibns\_20u}
\AttributeTok{  }\FunctionTok{name}\KeywordTok{:}\AttributeTok{ Output current}
\AttributeTok{  }\FunctionTok{min}\KeywordTok{:}\AttributeTok{ {-}5\%}
\AttributeTok{  }\FunctionTok{typ}\KeywordTok{:}\AttributeTok{ }\DecValTok{20}
\AttributeTok{  }\FunctionTok{max}\KeywordTok{:}\AttributeTok{ 5\%}
\AttributeTok{  }\FunctionTok{scale}\KeywordTok{:}\AttributeTok{ }\FloatTok{1e6}
\AttributeTok{  }\FunctionTok{digits}\KeywordTok{:}\AttributeTok{ }\DecValTok{3}
\AttributeTok{  }\FunctionTok{unit}\KeywordTok{:}\AttributeTok{ uA}

\FunctionTok{vgs}\KeywordTok{:}
\AttributeTok{  }\FunctionTok{src}\KeywordTok{:}
\AttributeTok{    }\KeywordTok{{-}}\AttributeTok{ vgs\_m1}
\AttributeTok{  }\FunctionTok{name}\KeywordTok{:}\AttributeTok{ Gate{-}Source voltage}
\AttributeTok{  }\FunctionTok{typ}\KeywordTok{:}\AttributeTok{ }\FloatTok{0.6}
\AttributeTok{  }\FunctionTok{min}\KeywordTok{:}\AttributeTok{ }\FloatTok{0.3}
\AttributeTok{  }\FunctionTok{max}\KeywordTok{:}\AttributeTok{ }\FloatTok{0.8}
\AttributeTok{  }\FunctionTok{scale}\KeywordTok{:}\AttributeTok{ }\DecValTok{1}
\AttributeTok{  }\FunctionTok{digits}\KeywordTok{:}\AttributeTok{ }\DecValTok{3}
\AttributeTok{  }\FunctionTok{unit}\KeywordTok{:}\AttributeTok{ V}
\end{Highlighting}
\end{Shaded}

Re-run the measurement and result generation

\begin{Shaded}
\begin{Highlighting}[]
\FunctionTok{make}\NormalTok{ typical OPT=}\StringTok{"{-}{-}no{-}run"}
\end{Highlighting}
\end{Shaded}

Open \texttt{results/tran\_Sch\_typical.html}

\subsection{Check waveforms}\label{check-waveforms}\index{check waveforms@Check waveforms}

You can either use ngspice, or you can use cicsim, or you can use
something I don't know about

Open the raw file with

\begin{Shaded}
\begin{Highlighting}[]
\ExtensionTok{cicsim}\NormalTok{ wave output\_tran/tran\_SchGtKttTtVt.raw }
\end{Highlighting}
\end{Shaded}

Load the results, and try to look at the plots. There might not be that
much interesting happening

\subsubsection{Searching waveforms}\label{searching-waveforms}

On the left side of the window you'll see a text box in the middle
between the filename, and the wave names. This is a regex search field,
and you can easily search for waveforms (like \texttt{i(v0)}) that you
want to find.

Note that the search field uses
\href{https://en.wikipedia.org/wiki/Regular_expression}{regular
expressions}. If you don't know regex, then it's time to learn. I always
use the perl regular expression variants.

For example, searching for ``i(v0)'' won't actually show anything,
because the \texttt{()} are special characters. ``i(v0)'' will find it
though.

I could search for both ibps and v0 at the same time with
\texttt{ibps\textbar{}i\textbackslash{}(}, so it's well worth learning.

A great resource is \href{https://regex.info/book.html}{Mastering
Regular Expressions}

\section{\texorpdfstring{All corners SPICE simulations
}{All corners SPICE simulations }}\label{all-corners-spice-simulations}

Analog circuits must be simulated for all physical conditions, we call
them corners. We must check high and low temperature, high and low
voltage, all process corners, and device-to-device mismatch.

\subsection{Remove Vh and Vl corners
(Makefile)}\label{remove-vh-and-vl-corners-makefile}

For the current mirror we don't need to vary voltage, since we don't
have a VDD.

Open Makefile in your favorite text editor.

Change all instances of ``Vt,Vl,Vh'' and ``Vl,Vh'' to Vt

\subsection{Run all corners}\label{run-all-corners}\index{run all corners@Run all corners}

To simulate all corners do

\begin{Shaded}
\begin{Highlighting}[]
\FunctionTok{make}\NormalTok{ typical etc mc}
\end{Highlighting}
\end{Shaded}

where etc is extreme test condition and mc is monte-carlo.

Wait for simulations to complete.

\subsection{Get creative with python}\label{get-creative-with-python}\index{get creative with python@Get creative with python}

Open \texttt{tran.py} in your favorite editor, try to read and
understand it.

The \texttt{name} parameter is the corner currently running, for example
\texttt{tran\_SchGtAmcttTtVt}.

The measured outputs from ngspice will be added to
\texttt{tran\_SchGtAmcttTtVt.yaml}

Delete the ``return'' line.

Add the following lines (they automatically plot the current and gate
voltage)

\begin{Shaded}
\begin{Highlighting}[]
\ImportTok{import}\NormalTok{ cicsim }\ImportTok{as}\NormalTok{ cs}
\NormalTok{fname }\OperatorTok{=}\NormalTok{ name }\OperatorTok{+}\StringTok{".png"}
\BuiltInTok{print}\NormalTok{(}\SpecialStringTok{f"Saving }\SpecialCharTok{\{}\NormalTok{fname}\SpecialCharTok{\}}\SpecialStringTok{"}\NormalTok{)}
\NormalTok{cs.rawplot(name }\OperatorTok{+} \StringTok{".raw"}\NormalTok{,}\StringTok{"time"}\NormalTok{,}\StringTok{"v(ibps\_5u),i(v0)"} \OperatorTok{\textbackslash{}}
\NormalTok{  ,ptype}\OperatorTok{=}\StringTok{""}\NormalTok{,fname}\OperatorTok{=}\NormalTok{fname)}
\end{Highlighting}
\end{Shaded}

Re-run measurements to check the python code

\begin{Shaded}
\begin{Highlighting}[]
\FunctionTok{make}\NormalTok{ typical etc mc OPT=}\StringTok{"{-}{-}no{-}run"}
\end{Highlighting}
\end{Shaded}

You'll see that cicsim writes all the png's. Check with
\texttt{ls\ -l\ output\_tran/*.png}.

You'll also notice it will slow down the simulation, so maybe remove the
lines from \texttt{tran.py} again ;-)

\subsection{Generate simulation
summary}\label{generate-simulation-summary}

Run

\begin{Shaded}
\begin{Highlighting}[]
\FunctionTok{make}\NormalTok{ summary}
\end{Highlighting}
\end{Shaded}

Install \href{https://pandoc.org}{pandoc} if you don't have it

Run

\begin{Shaded}
\begin{Highlighting}[]
\ExtensionTok{pandoc} \AttributeTok{{-}s}\NormalTok{   README.md }\AttributeTok{{-}o}\NormalTok{ README.html}
\end{Highlighting}
\end{Shaded}

to generate a HTML slideshow that you can open in browser. Open the HTML
file.

\subsection{Viewing results without GUI
browser}\label{viewing-results-without-gui-browser}

If you're on a system without a browser, or indeed a GUI, then it's
possible to view the results in the terminal.

Check if \texttt{lynx} is installed, if it's not installed, then

On linux

\begin{Shaded}
\begin{Highlighting}[]
\FunctionTok{sudo}\NormalTok{ apt{-}get install lynx}
\end{Highlighting}
\end{Shaded}

On Mac

\begin{Shaded}
\begin{Highlighting}[]
\ExtensionTok{brew}\NormalTok{ install lynx}
\end{Highlighting}
\end{Shaded}

Then

\begin{Shaded}
\begin{Highlighting}[]
\FunctionTok{lynx}\NormalTok{ README.html}
\end{Highlighting}
\end{Shaded}

\subsection{Think about the results}\label{think-about-the-results}\index{think about the results@Think about the results}

From the corner and mismatch simulation, we can observe a few things.

\begin{itemize}
\tightlist
\item
  The typical value is not 20 uA. This is likely because we have a M2
  VDS of 1 V, which is not the same as the VDS of M1. As such, the
  current will not be the same.
\item
  The statistics from 30 corners show that when we add or subtract 3
  standard deviation from the mean, the resulting current is outside our
  specification of +- 5 \%.
\end{itemize}

\section{\texorpdfstring{Draw Layout }{Draw Layout }}\label{draw-layout}

A foundry (the factory that makes integrated circuits) needs to know how
we want them to create our circuit. So we need to provide them with a
``layout'', the recipe, or instruction, for how to make the circuit.
Although the layout contains the same components as the schematic, the
layout contains the physical locations, and how to actually instruct the
foundry on how to make the transistors we want.

Open Magic VLSI

\begin{Shaded}
\begin{Highlighting}[]
\BuiltInTok{cd}\NormalTok{ work}
\ExtensionTok{magic}\NormalTok{ ../design/LELO\_EX\_SKY130A/LELO\_EX.mag}
\end{Highlighting}
\end{Shaded}

Now brace yourself, Magic VLSI was created in the 1980's. For its time
it was extremely modern, however, today it seems dated. However, it is
free, so we use it.

\subsection{Magic VLSI}\label{magic-vlsi}\index{magic vlsi@Magic VLSI}

Try google for most questions, and there are youtube videos that give an
intro.

\begin{itemize}
\tightlist
\item
  \href{https://www.youtube.com/watch?v=ORw5OaY33A4&t=9s}{Magic Tutorial
  1}
\item
  \href{https://www.youtube.com/watch?v=NUahmUtY814}{Magic Tutorial 2}
\item
  \href{https://www.youtube.com/watch?v=OKWM1D0_fPI}{Magic Tutorial 3}
\item
  \href{http://opencircuitdesign.com/magic/commandref/commands.html}{Magic
  command reference}
\item
  \href{https://analogicus.com/magic/}{Magic Documentation}
\end{itemize}

Default magic start with the BOX tool. Mouse left-click to select bottom
corner, left-click to select top corner.

Press ``space'' to select another tool (WIRING, NETLIST, PICK).

Type ``macro help'' in the command window to see all shortcuts

\begin{longtable}[]{@{}ll@{}}
\toprule\noalign{}
Hotkey & Function \\
\midrule\noalign{}
\endhead
\bottomrule\noalign{}
\endlastfoot
v & View all \\
shift-z & zoom out \\
z & zoom in \\
x & look inside box (expand) \\
shift-x & don't look inside box (unexpand) \\
u & undo \\
d & delete \\
s & select \\
Shift-Up & Move cell up \\
Shift-Down & Move cell down \\
Shift-Left & Move cell left \\
Shift-Right & Move cell right \\
\end{longtable}

\subsection{Add transistors}\label{add-transistors-1}\index{add transistors@Add transistors}

Open Cell -\textgreater{} Place Instance. Navigate to the right
transistor.

Place it. Hover over the transistor and select it with `s'. Now comes a
bit of tedious thing. Select again, and copy. It's possible to align the
transistors on-top of eachother, but it's a bit finicky.

Place all transistors on top of each other as shown below in the
picture.


{
\centering
\aicfig{media/LELO_EX_place.png}
}

\emph{Figure 2: Magic with the transistor instances stacked on top of
each other, shown as boxes on the left and with layers drawn on the
right}

\subsection{Place devices}\label{place-devices}\index{place devices@Place devices}

You will find that one of the more time consuming things with analog
layout is to place the devices, and to follow the design rules from
foundry. I detest tedious work. As such, I've tried for the past 25
years to simplify analog layout. I've not finished yet, but maybe you'll
find some of the scripts useful.

Note that the command below will override all your hard work ;-)

\begin{verbatim}
cd work
make xsch
cicpy sch2mag LELO_EX_SKY130A LELO_EX
\end{verbatim}

\subsection{Add Ground}\label{add-ground}\index{add ground@Add ground}

In the command window, type

\begin{Shaded}
\begin{Highlighting}[]
\OtherTok{see}\NormalTok{ no *}
\OtherTok{see}\NormalTok{ viali}
\OtherTok{see}\NormalTok{ locali}
\OtherTok{see}\NormalTok{ m1}
\OtherTok{see}\NormalTok{ via1}
\OtherTok{see}\NormalTok{ m2}
\end{Highlighting}
\end{Shaded}

Make a box around the layout by left clicking bottom left, and right
clicking top right. Press `x' to expand.

Change grid to 1 um. Set ``Window-\textgreater Snap to grid on''

Select a 1 um box below the transistors and paint the rectangle with
locali (middle click on locali)

Change to the `wire tool' with spacebar. Set ``Window-\textgreater{}
Snap to grid off''

Connect guard rings to ground.

Press the top transistor `S' and draw all the way down to connect all of
the transistors' source terminals. Use `shift-right click' to change
layer down


{
\centering
\aicfig{media/ground.png}
}

\emph{Figure 3: The source terminals of all transistors connected down
to the ground rail, DRC clean}

\subsection{Route Gates}\label{route-gates}\index{route gates@Route gates}

Press ``space'' to enter wire mode. Left click on the top gate to start
a wire, and right click to end the wire.

The drain of M1 transistor needs a connection from gate to drain. We do
that for the middle transistor. Change to the box tool (spacebar a few
times). Create a box that matches the locali. Connect the drain to the
gate in locali.


{
\centering
\aicfig{media/gates.png}
}

\emph{Figure 4: The gates routed together, with the gate to drain
connection of M1 made in locali}

\subsection{Drain of M2}\label{drain-of-m2}\index{drain of m2@Drain of M2}

Use the wire tool to draw connections for the drains.

To add vias you can do ``shift-left click'' to move up a metal, and
``shift-right click'' to go down.

It's a very good idea to have direction rules for metal layers. I would
recommend that you route metal1 vertical, metal2 horizontal, metal3
vertical etc. For locali it's usually all over the place.


{
\centering
\aicfig{media/drains.png}
}

\emph{Figure 5: The drain connections routed with the wire tool, using
vias to change metal layer}

\subsection{Add labels}\label{add-labels}\index{add labels@Add labels}

All ports must be named (IBPS\_5U, IBNS\_20U, VSS). The cicpy script may
add ports, but not necessarily where you want them.

Select a box on a metal, and use ``Edit-\textgreater Text'' to add
labels for the ports. Select the port button.

\section{\texorpdfstring{Layout verification
}{Layout verification }}\label{layout-verification}

The DRC can be seen directly in Magic VLSI as you draw.

To check layout versus schematic navigate to work/ and do

\begin{Shaded}
\begin{Highlighting}[]
\NormalTok{make cdl lvs}
\end{Highlighting}
\end{Shaded}

Remember to save the layout first.

If you've routed correctly, then the LVS should be correct.


{
\centering
\aicfig{media/layout.png}
}

\emph{Figure 6: The completed LELO\_EX layout with labelled ports, ready
for the LVS check against the schematic}

\section{\texorpdfstring{Extract layout parasitics
}{Extract layout parasitics }}\label{extract-layout-parasitics}

With the layout complete, we can extract parasitic capacitance.

\begin{Shaded}
\begin{Highlighting}[]
\FunctionTok{make}\NormalTok{ lpe}
\end{Highlighting}
\end{Shaded}

Check the generated netlist

\begin{Shaded}
\begin{Highlighting}[]
\FunctionTok{cat}\NormalTok{ lpe/LELO\_EX\_lpe.spi}
\end{Highlighting}
\end{Shaded}

\section{\texorpdfstring{Simulate with layout parasitics
}{Simulate with layout parasitics }}\label{simulate-with-layout-parasitics}

Navigate to sim/LELO\_EX. We now want to simulate the layout.

The default \texttt{tran.spi} should already have support for that.

Open the Makefile, and change

\begin{Shaded}
\begin{Highlighting}[]
\VariableTok{VIEW}\OperatorTok{=}\NormalTok{Sch}
\end{Highlighting}
\end{Shaded}

to

\begin{Shaded}
\begin{Highlighting}[]
\VariableTok{VIEW}\OperatorTok{=}\NormalTok{Lay}
\end{Highlighting}
\end{Shaded}

\subsection{Typical simuation}\label{typical-simuation}\index{typical simuation@Typical simuation}

Run

\begin{Shaded}
\begin{Highlighting}[]
\FunctionTok{make}\NormalTok{ typical}
\end{Highlighting}
\end{Shaded}

\subsection{Corners}\label{corners}\index{corners@Corners}

Navigate to sim/LELO\_EX. Run all corners again

\begin{Shaded}
\begin{Highlighting}[]
\FunctionTok{make}\NormalTok{ all}
\end{Highlighting}
\end{Shaded}

\subsection{Simulation summary}\label{simulation-summary}\index{simulation summary@Simulation summary}

Open \texttt{summary.yaml} and add the layout files.

\begin{Shaded}
\begin{Highlighting}[]
\AttributeTok{      }\KeywordTok{{-}}\AttributeTok{ }\FunctionTok{name}\KeywordTok{:}\AttributeTok{ Lay\_typ}
\AttributeTok{        }\FunctionTok{src}\KeywordTok{:}\AttributeTok{ results/tran\_Lay\_typical}
\AttributeTok{        }\FunctionTok{method}\KeywordTok{:}\AttributeTok{ typical}
\AttributeTok{      }\KeywordTok{{-}}\AttributeTok{ }\FunctionTok{name}\KeywordTok{:}\AttributeTok{ Lay\_etc}
\AttributeTok{        }\FunctionTok{src}\KeywordTok{:}\AttributeTok{ results/tran\_Lay\_etc}
\AttributeTok{        }\FunctionTok{method}\KeywordTok{:}\AttributeTok{ minmax}
\AttributeTok{      }\KeywordTok{{-}}\AttributeTok{ }\FunctionTok{name}\KeywordTok{:}\AttributeTok{ Lay\_3std}
\AttributeTok{        }\FunctionTok{src}\KeywordTok{:}\AttributeTok{ results/tran\_Lay\_mc}
\AttributeTok{        }\FunctionTok{method}\KeywordTok{:}\AttributeTok{ 3std}
\end{Highlighting}
\end{Shaded}

Run summary again

\begin{Shaded}
\begin{Highlighting}[]
\FunctionTok{make}\NormalTok{ summary}
\ExtensionTok{pandoc} \AttributeTok{{-}s}\NormalTok{  README.md }\AttributeTok{{-}o}\NormalTok{ README.html}
\end{Highlighting}
\end{Shaded}

Open the README.html and have a look a the results. The layout should be
close to the schematic simulation.

\section{Make documentation}\label{make-documentation}\index{make documentation@Make documentation}

Make a file (or it may exists)
\texttt{design/LELO\_EX\_SKY130A/LELO\_EX.md} and add some documentation
of what you've made.

Add the simulation results to your git repository to keep track

\begin{verbatim}
git add sim/LELO_EX/results/*.html
git add sim/LELO_EX/README.md 
\end{verbatim}

\section{Edit info.yaml}\label{edit-info.yaml}\index{edit info.yaml@Edit info.yaml}

Finally, let's setup the \texttt{info.yaml} so that all the github
workflows run correctly.

Mine will look like this.

You need to setup the url (probably something like
\texttt{\textless{}your\ username\textgreater{}.github.io}) to what is
correct for you.

I've added the doc section such that the workflows will generate the
docs.

The sim is to run a typical simulation.

\begin{Shaded}
\begin{Highlighting}[]
\FunctionTok{library}\KeywordTok{:}\AttributeTok{ LELO\_EX\_SKY130A}
\FunctionTok{cell}\KeywordTok{:}\AttributeTok{ LELO\_EX}
\FunctionTok{author}\KeywordTok{:}\AttributeTok{ Carsten Wulff}
\FunctionTok{github}\KeywordTok{:}\AttributeTok{ wulffern}
\FunctionTok{tagline}\KeywordTok{:}\AttributeTok{ The answer is 42}
\FunctionTok{email}\KeywordTok{:}\AttributeTok{ carsten@wulff.no}
\FunctionTok{url}\KeywordTok{:}\AttributeTok{ wulffern.github.io}
\FunctionTok{doc}\KeywordTok{:}
\AttributeTok{  }\FunctionTok{libraries}\KeywordTok{:}
\AttributeTok{    }\FunctionTok{LELO\_EX\_SKY130A}\KeywordTok{:}
\AttributeTok{      }\KeywordTok{{-}}\AttributeTok{ LELO\_EX}
\end{Highlighting}
\end{Shaded}

\section{Setup github pages}\label{setup-github-pages}\index{setup github pages@Setup github pages}

Go to your GitHub repository (repo). Press Settings. Press Pages. Choose
source under Build and Deployment -\textgreater{} GitHub Actions

Wait for the workflows to build. And check your github pages. Mine is
\url{https://wulffern.github.io/lelo_ex0_sky130a/}.

\section{Frequently asked questions}\label{frequently-asked-questions}\index{frequently asked questions@Frequently asked questions}

\emph{Q:} My GDS/LVS/DRC action fails, even though it works locally.

Sometimes the reference to the transistors in the magic file might be
wrong. Open the .mag file in a text editor and check. The correct way is

\begin{Shaded}
\begin{Highlighting}[]
\ExtensionTok{use}\NormalTok{ JNWATR\_NCH\_4C5F0  JNWATR\_NCH\_4C5F0\_0 ../LELO\_ATR\_SKY130A}
\end{Highlighting}
\end{Shaded}

It's the last \texttt{../LELO\_ATR\_SKY130A} that sometimes is missing.

\section{Summary}\label{summary}

The one-page version of this chapter:

\begin{itemize}
\tightlist
\item
  One cell end to end: schematic in xschem, a symbol, a testbench, and a
  cicsim run with a results table
\item
  The names are not cosmetic: instances, ports and directories follow
  the aicex conventions the tools expect
\item
  By the end you have run corners from one command - the pattern every
  later week repeats
\end{itemize}
